\documentclass[../main.tex]{subfiles}

\begin{document}

\chapter{Ellis group}
\label{ch:Ellis}

In this chapter we evaluate the Ellis group for some theories. Note in particular that the evaluation does not depend on the chosen model, as theorem ... holds.
The typical setting is the following:
\begin{itemize}
	\item	$L$ is a f.o.l. %finite?
	\item	$T$ is a complete theory for $L$ with Q.E. and countable models.
	\item	$M$ is a model of $T$. We want to evaluate $E_M$, the associated Ellis group.
	\item	$N$ is a $|M|^+$-saturated elementary extension of $M$.
	\item	$x,y$ are countable tuples of variables.
	\item	$\gamma \subseteq L$ is a set of formulas of the form $\phi(x;y)$, closed under boolean combinations.
		We also assume $L_y \subseteq \gamma$ 
		(This request is needed to have $p(y) \supseteq tp(m/\emptyset)$ in duality theorem). 
		%\gamma should be finite?
		%This is possible because $|L| < \omega$ and $T$ has Q.E.

	\item	We will denote, for every $n \in N^x$, $(n)_{\gamma,M} := \{a \in N^x : a \equiv_{\gamma,M} n\}$. Note that $N/\equiv_{\gamma,M} \cong S_x(M)$.
		Since $\gamma$ will always be fixed, at some point we will just omit it.
	\item	$\alpha(y)$ is a finitary type, meaning: a parameter free, complete type in which only a finite number of variables from $y$ actually appear.
	\item	Given $\phi_1(x;y),..,\phi_n(x;y) \in \gamma$ and a finitary type $\alpha(y)$, the couple $t = (\phi_1(x;y),..,\phi_n(x;y);\alpha(y))$ is called a pattern. 
	\item	$\LLang$ is the extended pattern language as described in Appendix A, that is: $\LLang = \{\R_t : t \text{ is a pattern}\}$.
		Recall that $S_\gamma(M)$ is the $\LLang$ structure whose domain is $N^x / \equiv_{\gamma,M}$ and such that 
		\[
			S_\gamma(M) \models \R_t((a_1)_{\gamma,M},..,(a_n)_{\gamma,M}) \iff \neg \exists b \in \alpha(M^y) (N \models \bigwedge_{i=1,..,n} \phi_i(a_i,b))
		\]
	\item	$\JJ$ is the associated universal e.c. model.
	\item	$\U$ is a monster model for $T$.
	\item	$\nonfork$ will always mean strong coheir (see --definition below--). This is needed to have --duality--.
	\item	Given $a,b$ tuples in $\U$ we will write $a \equiv y \nonfork_M b$ for the conjunction of the types $a \equiv y$ and $y \nonfork_M b$.
		Consequently $a \equiv \U^y \nonfork_M b$ stands for the set of realizations in $\U$. Informally we call such realizations independent sequences.
\end{itemize}

The strategy used for the evaluation relies on the following results:
\begin{enumerate}
	\item	--theorem section below-- gives a duality between $\HSS$ and independent sequences. 
		In most of the cases it is easier to determine the latter.
	\item	At this point assume we know $Hom(S,S)$. We look for the most non-injective morphisms.
		Let $g \in \HSS$ be such a morphism and let $f: \mathcal{\JJ} \to S$ be an embedding (such $f$ exists by theorem ..).
		Then $g \circ f: \JJ \to S$ must be an embedding (by theorem ..) so in particular must be injective.
		Let $p,q \in S$ such that $g(p) = g(q)$. Then it can't be $p,q \in ran(f)$.
	\item	From the previous point one has some incompatibilities between types.
		Combined with theorem .. stating that only maximal types survive in $\mathcal{J}$, usually this is enough to compute $\JJ$.
	\item	From theorem .. we know $E_M \cong Aut(\JJ)$.
\end{enumerate}

For a first example of this strategy see --link example empty language-- which can be considered in some sense a tutorial.

\section{Duality}
\label{sec:duality}

Enunciare e dimostrare dualità con nozione di strong coheir.


--duality Hrushovski-- establishes a duality between morphisms and coheirs.
In this section we prove that the same result holds for the richer version of the pattern language if we replace coheirs with finitary coheirs.

\begin{definition}
	We say that $tp(c/N)$ is a $\gamma$-finitary coheir of $tp(c/M)$ if for every $\phi(x;y) \in \gamma, n \in N^x, \alpha(y)$ finitary type we have
	\[
		c \models \phi(n,y) \land \alpha(y) \Rightarrow m \models \phi(n,y) \land \alpha(y) 
	\]
	We write $c \nonfork_{\gamma, M} N$.
\end{definition}

Note that finitary coheir implies coheir when we restrict to $\gamma$, as $N \supseteq M$.

\begin{theorem}

	Let $B$ countable model of $T$ and let $b$ enumeration of $B$.
	Then there exists a one to one correspondence between $H := \{h: S_{\gamma}(M) \to S_{\gamma}(B) : h \ \LLang-morphism\}$ and $C := (b \equiv \U \nonfork_M N)/\equiv_{\phi,M}$.
	The correspondence is given by:
	\begin{align*}
		f: H \to C, h \mapsto (c)_{\gamma,M} & \text{ with }c \models \phi(m,y) \iff h(m) \models \phi(x,b) \\
		g: C \to H, (c)_{\gamma,M} \mapsto h & \text{ with }h((n)_{\gamma,M}) \models \phi(x,b) \iff c \models \phi(n,y)
	\end{align*}

\end{theorem}

\begin{example}
	esempio linguaggio vuoto, teoria degli insiemi infiniti
\end{example}

\begin{example}
	esempio linguaggio <, teoria dlo
\end{example}

\begin{example}
	esempio teoria acf
\end{example}


\end{document}
